% Generated by scripts/generate-worked-calculations.R; do not edit by hand.
% Registry ID: est.mean.two_sided_finite_95
% Source notebook: notebooks/support/population-estimation/support.Rmd
Using the values from \emph{Case: Salaries}, the two-sided 95\% confidence interval\index{confidence interval}\index{interval!confidence} is now calculated as

\begin{equation*}
3,014.16 \pm 2.010 \cdot \frac{1,537.55}{\sqrt{50}}\sqrt{\frac{2,222 - 50}{2,222-1}} = 3,014.16 \pm 432.12
\end{equation*}

By considering the finite-population correction factor, the precision achieved\index{precision!achieved} is reduced from 436.97 to 432.12.

It is always permissible to use the finite-population correction factor regardless of the value of $n / N$, but if the sample is less than 10\% of the population, then the correction is close to 1 (no correction). In our case, the sample proportion was 50 / 2,222 $\approx$ 2.2\%, and the correction $\sqrt((2,222 - 50)/(2,222-1))$ $\approx$ 0.989.
